\chapter[Petroleum, Natural Gas, and Fuels]{Petroleum, Natural\protect\\ Gas, and Fuels}

\begin{kaichang}
Two flames burn in the evening kitchen. The gas stove has a quiet, clean-blue flame, barely luminous, making only the heated air shimmer. A leftover birthday candle at the table's corner burns bright yellow, flickering with a nearly invisible thread of black smoke at its tip.

Look at a lighter too: its clear case holds half a tank of colorless liquid that sloshes like water, yet pressing it releases ignitable gas.

Guess three things: how do the blue and yellow flames' fuels differ? Why does the lighter's liquid become gas on emerging? And if natural gas is odorless, where does the sharp stink of a household gas leak come from?

Write down your guesses. By the end, these three questions will form one connected answer.
\end{kaichang}

\section{First, Look Closely at Flames}

\textbf{Translator safety note:} Do not adjust a stove's air shutter, deliberately make a yellow flame or soot, or burn or smell fuel for these observations. Incomplete combustion can produce deadly carbon monoxide; flame color and smell do not establish safety. Leave appliance adjustment to a qualified technician and use a vetted demonstration video. If a CO alarm sounds, go outside and call emergency services. See \href{https://www.cpsc.gov/safety-education/safety-guides/carbon-monoxide/carbon-monoxide-fact-sheet}{CPSC carbon-monoxide guidance}.

Before entering molecules, use eyes, nose, and a cold plate to examine burning itself.

Hold an ice-cold ceramic plate a dozen or so centimeters above a gas-stove flame, taking care against burns. Water droplets soon cover its underside: the flame makes water. Bring the same plate above a candle, and besides droplets it gains black soot that leaves a dark streak under a finger. Blue flames leave almost none; yellow flames do.

Observe shapes too. A candle has a small dark-blue region at its base, a bright-yellow main body, and a darkening tip; stove flames form neat blue cones. Reduce the stove's air shutter opening to admit less air, and the flames gradually lengthen and turn yellow, blackening pan bottoms. This is how older gas stoves soot up cookware.

Your nose offers clues too. Pure gasoline and alcohol produce hardly any smoke smell while burning; the white plume when a candle is blown out smells acrid; diesel exhaust is black and choking. All are combustion, but different products reveal themselves in color, smell, and soot.

\textbf{Translator safety note:} For the candle activity below, an adult must control the flame and hot objects. Do not rub soot with your fingers; observe it without contact, allow the spoon to cool completely, and wash hands afterward. Never burn gasoline or alcohol as an extension of this activity. See \href{https://www.acs.org/education/celebrating-chemistry-editions/2024-ccew/science-safety-tips.html}{ACS activity precautions} and \href{https://www.poison.org/articles/rubbing-alcohol-only-looks-like-water}{Poison Control's alcohol-fire warning}.

\begin{shiyan}
\textbf{Take a Flame's Fingerprint}

Materials: a candle; a metal spoon or ceramic dish; matches; a cup of water; an accompanying adult.

Procedure:
\begin{enumerate}[itemsep=1pt]
\item Ask an adult to light the candle on a table away from curtains and paper.
\item Hold the spoon handle, place its back in the upper outer flame for two or three seconds, then withdraw immediately.
\item Observe the back: first a watery mist, then black powder visible after wiping dry. Rub a little between your fingers.
\item Repeat above the flame's top, farther from its core, and compare the soot amounts.
\end{enumerate}

Safety: heat for only two or three seconds. The spoon's back becomes hot: do not touch it. Blow out the candle immediately afterward; an adult must remain present throughout. The powder is almost pure carbon, evidence that the candle has not burned completely. Remember it: we will account for it later.
\end{shiyan}

\section{Fuel's Identity: Carbon and Hydrogen Chains}

What burns in these flames? Microscopically, the answer is surprisingly repetitive: natural gas, lighter gas, gasoline, kerosene, diesel, candles, and asphalt all consist mainly of \textbf{alkanes}.

Chapter 36 gave carbon four hands: each atom forms four single bonds, joining chains and rings. Alkanes are the simplest family, with only single carbon--carbon bonds and every remaining hand filled by hydrogen. One carbon and four hydrogens make methane, \chem{CH_4}, natural gas's main component. Two carbons and six hydrogens make ethane, \chem{C_2H_6}; three make propane, four butane. The lighter's water-like liquid is pressurized liquefied butane, \chem{C_4H_{10}}. Each chain extension adds \chem{CH_2}, like building blocks continuing through pentane, hexane, heptane, and chains dozens of carbons long.

With every carbon hand occupied, alkanes are \textbf{saturated hydrocarbons}. Full occupancy means few interfaces: Chapter 37 placed reactions at functional groups, which alkanes nearly lack. They are mild-mannered, not casually reacting with acids or bases, not readily discoloring or deteriorating, and resting in tanks for years. They enthusiastically do just one thing: burn. Fuel wants exactly that temperament, quiet until ignition.

Two properties determine their fate: \textbf{very low molecular polarity, with only weak van der Waals attractions between molecules}, Chapter 22's weak handshakes. First, they are water-insoluble: oil and water always separate, and rainwater does not dissolve service-station oil slicks into drains but carries them floating. Second, boiling point depends almost entirely on chain length: shorter chains hold one another more loosely and become gas more readily.

\begin{table}[h]
\centering
\begin{tabular}{clll}
\toprule
Carbons & \shortstack[l]{Room-temperature\\state} & Boiling range & Everyday uses \\
\midrule
1--4 & Gas & \shortstack[l]{Below roughly\\$0\,^{\circ}\mathrm{C}$} & \shortstack[l]{Natural gas, lighter gas,\\liquefied petroleum gas} \\
5--12 & Volatile liquid & About $30$--$200\,^{\circ}\mathrm{C}$ & Gasoline, solvents \\
10--16 & Liquid & About $150$--$300\,^{\circ}\mathrm{C}$ & Kerosene, jet fuel \\
14--20 & Viscous liquid & About $250$--$350\,^{\circ}\mathrm{C}$ & Diesel \\
Above 20 & Semisolid to solid & Higher & \shortstack[l]{Lubricating oil,\\paraffin wax, asphalt} \\
\bottomrule
\end{tabular}
\caption{The alkane family sorted by chain length. Boiling ranges are approximate because branched isomers are also present.}
\end{table}

This table maps the chapter: natural gas is the shortest-chain end, gasoline the middle, candles and asphalt the longest. Not different raw materials, but one family lined up by boiling point.

Correct another intuition: petroleum is not only burned. At least a dozen things on or around you today come directly from it: synthetic clothes, plastic cups, sneaker soles, shampoo bottles, ointments, synthetic-rubber tires, road asphalt. Modern industry's more valuable use is as \textbf{chemical feedstock}: split and rearrange its chains to build plastics, fibers, drugs, and dyes. Chapter 41 treats dyes; Chapter 48, polymers such as plastics. Burning whole barrels of this molecular building-block warehouse for heat looks somewhat extravagant to chemists. That is another reason for energy transition: let sunlight and wind supply energy and leave carbon chains for materials.

\section{A Tower That Sorts by Boiling Point}

Crude oil pumped underground is a thick black stew of alkanes, cyclic hydrocarbons, and aromatics, from one carbon to more than forty, unusable directly. A refinery's first major task is not making gasoline but \textbf{sorting}, separating the mixture by boiling point in a fractional distillation column.

The principle is as plain as steaming buns. Chapter 22 described boiling as molecules escaping neighbors' handshakes; shorter chains grip weakly and boil lower. Heat crude above three hundred degrees Celsius and feed it into a tower, actually a steel column tens of meters high, hot below and cold above. Vapor rises until a level cools it below its boiling point, where it condenses into trays. Heat-resistant asphalt stays below; diesel and kerosene collect midway; gasoline climbs higher; light methane and ethane leave the top as gas. One tower sorts the stew onto a dozen or so shelves.

\begin{qiaoliang}
Why can different boiling points separate substances instead of leaving them mixed forever? Evaporation and condensation run both ways: molecules continually enter and leave a liquid surface at any temperature. Lower-boiling components always occupy a larger share of vapor than liquid, applying Chapter 31's equilibrium reasoning. Each tray effectively condenses and re-evaporates vapor; dozens of trays mean dozens of redistributions, screening low-boiling components upward. More levels give purer separation. No chemical bond breaks: the molecules entering and leaving are the same, merely sorted. \textbf{Fractional distillation is a physical change, not production of new substances.}
\end{qiaoliang}

Sorting faces a practical problem: demand differs among shelves. Crude naturally contains too many long chains and too little five-to-twelve-carbon gasoline for the world's cars. Refining's second task is therefore actual manufacture: \textbf{cracking} unsalable long molecules into shorter ones, and \textbf{reforming} short chains into more useful shapes.

Cracking uses heat, pressure, and catalysts. Chapter 30 said catalysts change routes and rates, not equilibrium positions. Break a sixteen-carbon diesel molecule at a carbon--carbon single bond to produce, for example, an eight-carbon alkane and an eight-carbon alkene. This is a true chemical change: bonds break and new molecules appear.

Reforming tackles \textbf{knock}. Gasoline engines use spark plugs, with flames ideally sweeping smoothly across cylinders. Too many straight-chain alkanes let the compressed mixture ignite by itself before the flame arrives, striking the cylinder with a knock, wasting fuel and damaging engines. People found that eight-carbon straight n-octane knocks readily, while multiply branched isooctane resists. Reforming bends straight chains into branched and cyclic structures, improving resistance.

Engineers quantify this with \textbf{octane rating}: isooctane is assigned 100 and n-heptane 0. Compare a gasoline's knock resistance with mixtures of these references; the equivalent isooctane percentage gives the rating. Station labels 92 and 95 refer to this scale, not 95\% purity or three extra energy units. Higher octane need not make a car more powerful; it simply resists premature self-ignition in high-pressure engines.

\section{Combustion: Settling the Energy Accounts}

Now the chapter's principal reaction. Natural gas burns in one line:

\[\chem{CH_4} + \chem{2O_2} \rightarrow \chem{CO_2} + \chem{2H_2O}\]

Complete butane combustion follows the same pattern with larger coefficients:

\[\chem{2C_4H_{10}} + \chem{13O_2} \rightarrow \chem{8CO_2} + \chem{10H_2O}\]

Only carbon dioxide and water appear. This is \textbf{complete combustion}: ample oxygen lets each carbon take two oxygens, and each hydrogen its oxygen to form water, burning cleanly. Blue stove flames and plate droplets show it happening.

Where does energy come from? Chapter 27 calculated it, but repeat the easily confused rule: \textbf{breaking bonds costs; forming bonds pays}. Igniting methane first absorbs heat to break \chem{C-H} and \chem{O=O} bonds, requiring a match or spark and Chapter 29's activation energy. Carbon and oxygen, hydrogen and oxygen then form \chem{CO_2} and \chem{H_2O}, releasing heat. Income far exceeds cost, so the net reaction heats, paying neighboring molecules' ignition fees and sustaining itself. Flame energy is not stored heat hidden in fuel: \textbf{product bonds have a cheaper energy account than reactant bonds}, returning the difference as light and heat.

What if oxygen is insufficient? Carbon cannot obtain enough and burns halfway, yielding two unfinished products: carbon monoxide, \chem{CO}, with one oxygen per carbon, and unburned carbon particles, the soot on your spoon. This is \textbf{incomplete combustion}: yellow candle flames, blackened pans, and black diesel exhaust all display it.

Yellow is an accidental light show. Carbon particles heat above a thousand degrees Celsius and glow like hot iron, as Chapter 8 explained for hot objects. Countless glowing specks make a bright-yellow flame. Blue flames have no glowing carbon particles, only faint molecular emission from reacting gases, producing quiet blue. The first opening answer is emerging.

Incomplete combustion does more than waste fuel. One unfinished product, colorless and odorless carbon monoxide, is a genuine killer deserving its own section.

\section{Carbon Monoxide: A Passenger Blocking the Station}

Your oxygen transport is precise logistics: lungs load cargo; hundreds of millions of hemoglobins in blood are trucks. At each truck's center is an iron-atom parking space fitting one oxygen molecule, \chem{O_2}, described in Chapter 15's coordination chemistry. It carries oxygen to tissues, unloads, and returns empty.

Unfortunately, carbon monoxide resembles oxygen in shape and charge distribution and fits the same space, gripping more than two hundred times more tightly. Inhaled carbon monoxide occupies truck after truck long-term, preventing oxygen loading and paralyzing transport. Brain and heart run short first: headache, dizziness, nausea, then unconsciousness, with severe cases fatal within tens of minutes. Most insidiously, sleeping victims receive not even headache's warning.

\begin{wuqu}
``Gas smells, so I can detect a leak; I can detect carbon monoxide poisoning early too.'' Dangerously wrong. Methane, propane, and butane, the main components of natural gas and LPG, are \textbf{all odorless}, as is colorless carbon monoxide. The household gas smell is a tiny sulfur-containing odorant, a thiol, \textbf{deliberately added} by suppliers, detectable at extremely low levels: an artificial smell alarm. But incomplete combustion's carbon monoxide carries no protective odorant, so you cannot smell it. Only three defenses are truly reliable: ventilation, no gas water heaters or charcoal heating in closed rooms, and a carbon monoxide alarm. Charcoal hotpot cooking in a winter room with shut windows gambles with your oxygen transport.
\end{wuqu}

Remember the firm rule: \textbf{no flame belongs in a sealed room}. Gas water heaters require exhaust ducts, charcoal barbecue belongs outdoors, and cars must not idle in garages with closed doors. This chapter offers no home trial for these: dangerous carbon monoxide experiments are video- or teacher-demonstration only.

\section{Where Does Petroleum Come From?}

\begin{zhengju}
Petroleum lies in rock hundreds to thousands of meters underground. How was its birth certificate read? Simply digging down cannot reveal its past.

In the 1930s, German chemist Alfred Treibs did detective work by isolating porphyrins from petroleum and oil shale. You need not memorize the name, but these frameworks belong to two key biological molecules: chlorophyll, plants' sunlight catcher, and heme, hemoglobin's iron-containing part just discussed. Treibs found petroleum porphyrins almost directly matching chlorophyll and heme frameworks, with changed central metals and time-trimmed side chains. Like identity-card fragments at an accident scene, petroleum held their molecular fossils.

Objections arose: might petroleum coincidentally synthesize similar structures? More evidence accumulated. Hundreds and thousands of biomarkers were found, molecular fossils made only by organisms, such as hopanes derived solely from bacterial membranes. Carbon isotopes also voted: Chapter 7 introduced carbon-12 and carbon-13, and biological enzymes prefer lighter carbon-12, making biological carbon isotopically light. Petroleum carries this biological fingerprint, distinctly different from inorganic-rock carbon.

The evidence chain now closes: petroleum mainly comes from ancient marine plankton and algae buried in sediment, slowly transformed over millions to hundreds of millions of years under oxygen-poor, hot, high-pressure conditions. Natural gas often accompanies it as the lightest breakdown fraction. Coal mainly comes from ancient land plants. Together they are \textbf{fossil fuels}, literally fuel made from fossils. Every liter of gasoline you burn is sunshine stored by plankton on a bright day hundreds of millions of years ago.
\end{zhengju}

\section{Burning's Cost: Carbon and the Greenhouse Effect}

Burning fossil fuels moves carbon buried for hundreds of millions of years back into air within centuries. Make a realistic estimate: a household car burns about a thousand liters of gasoline annually. Approximate it as octane, \chem{C_8H_{18}}, with density 0.75 kilograms per liter: 750 kilograms. Each molecule's eight carbons become eight \chem{CO_2}. Carbon's atomic mass is 12; \chem{CO_2}'s formula mass 44, so twelve parts carbon become forty-four parts carbon dioxide. About 85 percent of 750 kilograms, or 640 kilograms, is carbon. Multiply by $44/12$ to get about \textbf{2300 kilograms of carbon dioxide}: the car emits 2.3 tonnes of invisible gas annually, more than its own weight. Over a billion cars worldwide, plus power stations, ships, and planes, move tens of billions of tonnes of carbon into air through human activity each year.

Where does the carbon dioxide go, and what does it do? Chapter 55 fully explains the carbon cycle. For now, carbon constantly moves among atmosphere, ocean, and land life; fossil burning \textbf{adds an extra supply} from outside that cycle. Carbon dioxide transmits visible sunlight but absorbs outgoing terrestrial infrared radiation, like adding a blanket: the \textbf{greenhouse effect}. Two qualifications matter. First, the effect itself is beneficial; without it, Earth's average temperature would be below minus ten degrees. The problem is rapid addition: atmospheric \chem{CO_2} has risen from about 280 ppm before industrialization to over 420 ppm today, thickening the blanket too quickly for smooth climatic adjustment. Second, climate is highly complex, with clouds, currents, and ice feedbacks. Science gives probability ranges for particular places and years, yet the core conclusion that fossil burning raises global average temperature has evidence comparable to smoking's harm to health.

The answer is not to stop using energy but change carbon's sources and destinations. Wind, solar, and nuclear power do not move carbon; electric vehicles shift emissions from countless tailpipes to controllable power stations. Captured \chem{CO_2} can be returned underground or combined with green hydrogen to remake fuel and close the loop. Each \textbf{energy-transition} technology rests on earlier principles: batteries (Chapter 34), catalysts (Chapter 30), and energy accounting (Chapter 27). Chemistry did not create the problem, but solutions cannot bypass it.

\section{Boundaries: Where This Model Falls Short}

While fresh in mind, mark the model's limits:

\begin{itemize}[itemsep=2pt]
\item \textbf{Alkanes lining up by chain-length boiling point is simplified}. Real gasoline mixes hundreds of molecules; branched isomers differ from straight chains, so fractions have boiling ranges rather than single values.
\item \textbf{Octane rating is empirical, not a precise molecular property}. It matters under standard engine-test conditions; oxygenated fuels such as ethanol can exceed 100. It does not measure energy, so higher-grade gasoline is not more powerful.
\item \textbf{Complete combustion is an ideal limit}. Real flames always combine complete and incomplete burning in differing proportions; vehicle exhaust always contains carbon monoxide and particles, requiring a three-way catalytic converter, Chapter 30's exhaust chemical plant.
\item \textbf{The greenhouse effect is not burning's only consequence}. Particulates, nitrogen oxides, and sulfur oxides also create nearby pollution, affecting lungs rather than climate and requiring different control approaches.
\end{itemize}

\section{Back to the Kitchen Flames}

Now fulfill the opening promise.

\textbf{Blue and yellow flames:} stoves premix natural gas with air, supplying enough oxygen for clean methane burning without glowing carbon particles, leaving pale-blue molecular emission. Candle-wax vapor lacks oxygen near the core, forming particles that glow yellow when heated and finish burning at the outside. Thus candles divide their fire between illumination and combustion. The lighter's yellow flame mentioned at the start is similar: emerging butane lacks time to mix enough air.

\textbf{Why lighter liquid becomes gas:} butane boils around $-1\,^{\circ}\mathrm{C}$, so should be gaseous at ordinary pressure and temperature. Internal pressure holds it liquid, as Chapter 22's intermolecular handshakes explain. Opening the valve removes pressure, and it immediately boils outward as gas. The water-like liquid is gas temporarily imprisoned by pressure.

\textbf{The gas smell:} natural gas is odorless; added thiols are a supplier-installed chemical alarm. Carbon monoxide from incomplete combustion has no smell and requires ventilation and alarms.

Three flames, three questions, one logic: \textbf{molecular structure determines burning; completeness determines what remains}.

\section{Beyond Fuels: The Next Interface}

We met alkanes with almost no reaction interfaces, saving their whole interface budget for burning alone. Yet adding terminal hydroxyl, \chem{-OH}, from Chapter 37 changes everything: a molecule can dissolve in water, give scent, be recognized and dismantled by liver enzymes, and drive microbes wild. This little hydroxyl interface creates alcohol, vinegar, perfume, and disinfectant. Next chapter attaches our first functional group and explores its transformations.

\begin{tiaozhan}
Settle a combustion account yourself; skipping does not break the main thread. Approximate tabulated bond energies in $\mathrm{kJ/mol}$ are \chem{C-H}, 413; \chem{O=O}, 498; \chem{C=O} in \chem{CO_2}, 799; and \chem{O-H}, 463. For methane combustion, $\chem{CH_4} + \chem{2O_2} \rightarrow \chem{CO_2} + \chem{2H_2O}$, breaking costs $4\times 413 + 2\times 498 = 2648$, and forming pays $2\times 799 + 4\times 463 = 3450$. Net heat is about $3450-2648=802$ $\mathrm{kJ/mol}$, fairly close to the measured 890; the gap comes from average bond-energy approximations. Remember \textbf{cost first, income afterward}: ignition pays activation energy before the flame sustains itself. Why can carbon monoxide still burn and release heat? Going from \chem{CO} to \chem{CO_2} forms another bond, recovering the energy incomplete combustion wasted.
\end{tiaozhan}

\section*{Try It Yourself}

\begin{sikaoti}
\item Write and balance complete propane combustion, \chem{C_3H_8}, a principal LPG component. How many oxygen molecules does one propane consume? With half as much oxygen, what two carbon-containing products might form?
\item Sketch a fractional-distillation column: a vertical tower, bottom feed inlet, top-to-bottom temperature gradient, and outlets for natural gas, gasoline, kerosene, diesel, and asphalt at appropriate heights. Note which differing property allows separation.
\item Switching gasoline from 92 to 95 does not make a car more powerful. Explain what octane rating measures and which molecular structures give high ratings.
\item Use this chapter's estimate to calculate the carbon dioxide from a household's monthly use of twenty cubic meters of natural gas, treated as methane at 0.72 kilograms per cubic meter. Hint: \chem{CH_4} has formula mass 16, of which carbon contributes 12.
\item Someone heats a sealed winter room with charcoal, saying a door crack and absence of smell make it safe. Use particles to identify both errors.
\item The petroleum-porphyrin story says biological carbon is light. If a carbon molecule in Martian rock had a markedly light isotope ratio, would that prove life? What alternatives would need excluding?
\end{sikaoti}
